\documentclass[11pt]{article}
\usepackage[T1]{fontenc}
\usepackage[utf8]{inputenc}
\usepackage{lmodern,amsmath,amssymb,microtype,xcolor}
\usepackage[margin=.9in]{geometry}
\providecommand{\tightlist}{\setlength{\itemsep}{0pt}\setlength{\parskip}{0pt}}
\usepackage[colorlinks=true,urlcolor=blue!60!black]{hyperref}
\setlength{\parindent}{0pt}
\setlength{\parskip}{.5em}
\setlength{\emergencystretch}{2em}
\author{Wentao Li\\\small Prepared with assistance from OpenAI Codex}
\date{November 22, 2026}
\title{The number of colors matters for Ramsey finiteness on a fixed set}
\hypersetup{pdfauthor={Wentao Li},pdftitle={The number of colors matters for Ramsey finiteness on a fixed set}}
\begin{document}
\maketitle
Brot, Cao and Fernández-Bretón ask how their local finiteness classes change when one allows more than two colors. The theorem below separates every successive finite palette for ordinary Ramsey finiteness and for Boolean Hindman finiteness with two summands. The local qualifier matters: a global Ramsey theorem applying to every infinite set permits successive thinning, whereas a theorem about one set need not apply to its infinite subsets. \href{https://arxiv.org/abs/1910.11025}{2019 preprint, Question 5.1(5) in v1 and 5.1(4) in the final version} · \href{https://doi.org/10.1016/j.apal.2021.102961}{2021 journal article}

\textbf{Theorem.} Consistently with ZF, for each \(r,k\ge2\) there is an infinite \(X_{r,k}\) on which every \(k\)-coloring of \(r\)-element subsets has an infinite homogeneous set but some \((k+1)\)-coloring has none. For \(X_{2,k}\) the same strict threshold holds for colorings of finite subsets with the requirement that all one- and two-term symmetric differences of distinct members of an infinite family have one color. All these witnesses can be taken in one model.

\textbf{Finite structures.} Fix \(r,k\). Use two sorts, points and lines, with an incidence relation. Each \(r\)-set of points has one of \(k+1\) named colors. Each line has as its label a two-element subset of the palette, and every \(r\)-set lying on that line must have a color in its label. Any two distinct lines have at most \(r\) common points. A finite induced substructure is called closed in an extension if it already contains every common point of any two of its lines, and every line through any \(r+1\) of its points. There is at most one such line through those points. The closure conditions are transitive. A set of \(r\) points without lines is closed in every extension.

There is amalgamation over closed induced substructures. Identify the common base \(P\) and take the union of two finite extensions with no additional incidences. Two lines on opposite sides and outside the base can meet only in \(P\). If they shared \(r+1\) such points, both would already be in \(P\). A pair of old lines cannot acquire a new common point from the other side, since this would require both lines to be in \(P\). Nor can an old \((r+1)\)-set gain a new line. Now color each mixed \(r\)-set (one meeting both sides outside \(P\)). A line containing it must be in \(P\), and there is at most one: a point outside \(P\) cannot be common to two base lines. Hence there are at least two permissible choices for its color, independently of all the other mixed color assignments. This gives the required closed embeddings. Disjoint joint embeddings are included.

Take the countable homogeneous limit for these closed embeddings, constructed by successively realizing every finite closed extension requirement. There are only countably many finite isomorphism types and embeddings. Finite tuples are contained in finite closed substructures, every finite closed extension embeds over its base, and finite closed isomorphisms extend by back-and-forth. In particular the orbits of \(r\)-sets of points under the full automorphism group are exactly their colors. Lines of each label exist and each has infinitely many points: one can add any number of new points incident only to that line and color all new \(r\)-sets compatibly.

\textbf{The finite-support model.} For every \((r,k)\) take countably many individually named disjoint copies of its limit. Use the product of component automorphism groups on all the point and line atoms, and pointwise stabilizers of finite supports. Let \(X_{r,k}\) consist of the point atoms in the copies for the specified parameters. Given a supported coloring of \([X_{r,k}]^r\) by \(k\) colors, take a copy disjoint from its support. There the value depends only on the original color of the \(r\)-set. Two of the \(k+1\) original colors have the same value. A line labeled by those two colors gives an infinite supported homogeneous set of points.

For failure with \(k+1\) colors, use the original colors within a component and give a fixed color to sets spanning components. If \(Y\subseteq X_{r,k}\) is infinite and supported, enlarge the support to a finite closed base in each component it meets. Some \(y\in Y\) is outside their union. In its component take a finite closed extension of the local base containing \(y\). Freely copy this extension over the base \(r+1\) times, choosing no new incidences. Distinct copies of \(y\) share at most one base line, and all have exactly the same incidences with the base lines. Every \(r\)-set of these marked points can therefore be colored from a common palette of size at least two. Assign different colors on two such \(r\)-sets and embed the extension. Homogeneity puts every marked point into \(Y\), contradicting homogeneity for the original coloring.

\textbf{Two Boolean summands.} Take \(r=2\). For a \(k\)-coloring of finite subsets, the pair result gives an infinite set \(Z\) whose pairs all have one color. Fix \(z\in Z\) and take the pairs \(\{z,x\}\) for \(x\in Z\setminus\{z\}\). Symmetric differences of two distinct such pairs are again pairs from \(Z\). This is the source's Theorem 3.8 argument, with its palette retained.

For the other direction, call the original edge color \(d\) and regard its labels as distinct residues modulo \(k+1\), including on edges spanning components. Color a finite point set \(F\) by the sum of \(d(e)\) over all its edges, modulo \(k+1\). Let \(Y\) be any supported infinite family of finite point sets. Some \(F\in Y\) contains a point outside the finite closed bases for its support, since a finite union of bases has only finitely many subsets. In the finitely many components touched by \(F\) or the support, take finite closed extensions containing all of \(F\). Embed two copies over the bases. Their two marked finite sets \(F_1,F_2\) are in \(Y\) by componentwise homogeneity. Their symmetric difference includes the two copies of the chosen outside point. In its component the edge between those copies is mixed, with at least two allowed original colors. Vary just that one edge, leaving all other colors fixed, and embed each resulting finite structure. The two resulting symmetric differences have different colors modulo \(k+1\). This rules out every monochromatic infinite family in the required sense.

The joint assertion about the tagged family of witnesses transfers to pure ZF by Jech--Sochor. Tests consist of finite subsets, colorings, infinite homogeneous subsets or families, and pairs of their members. All are coded at uniformly bounded rank over the witnesses and \(\omega\). At arity one the finite pigeonhole principle applies in ZF. For full Boolean Hindman finiteness the source's characterization by Dedekind-finiteness of \([X]^{<\omega}\) already gives independence of the finite palette: in the positive case a countably infinite pairwise disjoint family gives the usual finite-color Hindman theorem on a countable subspace. The same observation covers the earlier collapses to the full property. Dependence of the \emph{three}-summand local property on the number of colors is a separate part of the original question and is not decided by these two-summand examples.

\textbf{Significance and priority.} Provisional significance \textbf{60/100} for two full finite-palette hierarchies in the same model. They address the ordinary Ramsey and two-summand Boolean clauses of an open-ended question. The incidence parameters let a supported homogeneous set select a pair of allowed colors, while closedness prevents the individual \(r\)-sets from naming that line. Independent proof and priority review are pending.
\end{document}
